\subsubsection{Lattice inclusions on the selected lift}\label{subsec:lift-exponent-models}
\hypertarget{source-lift-exponent-models}{}

The least exponents entering the signed index measure a further lattice inclusion beside the congruence length. To bound them, put the two coefficient lattices in the same space of forms at the selected lift.

Fix a selected lift, a tuple and a coefficient place \(w\mid p\).
Write \(Z\) for the selected space of forms at its row level.
In the coordinates of \(Z\), \(\mathcal C\) and \(\mathcal D\)
are the inverse images of the value lattices for the dual Weyl
and Weyl coefficients, respectively. The coordinate maps on
the finite class set identify both lattices simultaneously.
\begin{proposition}[Weyl inclusion on the original section lattices]
\label{prop:local-weyl-value-inclusion}
On the selected tuple and place \(w\mid p\), the sum of the exact Weyl exponents at all labels gives
\[
 p^{\sum_\tau e_p(\lambda_\tau)}\mathcal C_w\subseteq\mathcal D_w.
\]
For \(y<p\le h\) as in Lemma~\ref{hyp:weyl-prime-range},
exclude \(q,d,p_{\rm de},p'_{\rm de}\) and primes dividing an
integer of \(\mathcal I\). A power that gives the value inclusion at every local model under
an occurring integral Hecke character gives the same inclusion on
\(\mathcal C_w,\mathcal D_w\) at the original level.
\end{proposition}

\begin{proof}
\proofstep{Tensoring the exact inclusions and localizing the value map.}
At each coefficient label \(\tau\) choose a local basis of the
underlying lattice. The two coefficient models in that basis
use the dual Weyl and Weyl lattices \(W_\tau^\vee,W_\tau\) in
the prescribed rational realizations. The inclusion of
Section~\ref{subsec:lift-exponent-primes} tensors to
\[
 p^{\sum_\tau e_p(\lambda_\tau)}
          \bigotimes_\tau W_\tau^\vee
       \ \subseteq\  \bigotimes_\tau W_\tau .
\]
An additive weight translation twists both sides by the same
determinant line and does not change this exponent. If
\(\operatorname{ev}\) is the finite map of value coordinates,
the inverse image of the Weyl lattice is the kernel of its
composition with reduction modulo that lattice.
Writing \(q\) for reduction modulo the Weyl value lattice, exactness of localization gives
\[
 (\ker(q\circ\operatorname{ev}))_w=\ker(q_w\circ\operatorname{ev}_w).
\]
Equivalently one can
multiply the finitely many denominators outside \(w\) occurring
in the individual coordinates. The displayed inclusion therefore
gives the asserted exponent on each tuple.
\label{cor:lift-exponent-models-1}

\proofstep{Comparing the two lattices in the same operator space.}

We describe the sharper bound at primes in the weight range
\(y<p\le h\) outside the exceptional interval and specified
set, using the tuple assigned to \(w\). The base level and selected block hypotheses give finite free tensor value lattices over the same local coefficient ring. Choose
an embedding of the coefficient field into \(\overline{\Q}_p\)
inducing \(w\), and adjoin the roots of unity for the finite
level character to the induced local field. Its integer ring
\(O\) is a DVR. The integral frame data are changes of basis
on the two local value lattices, including the level group
matrices and character.

In the common interpolation frame a Hecke correspondence is represented on both lattices by the same operator
\[
 T=\sum_j a_j R(g_j),\qquad a_j\in O,
\]
with \(R(g_j)\) the specified right translations. Thus on the auxiliary space \(M\) of
integral forms the operators define an action of the commutative
polynomial algebra \(\mathcal T\) over \(O\) in these operators.
This action preserves the forms with values in the Weyl lattice,
whose submodule is \(M_{\rm W}\).

\proofstep{Detecting a failed inclusion at an occurring residual character.}
Take the integral form \(x_0\) from the dual Weyl reading of a
vector of \(\mathcal C\). Its value coordinates in the auxiliary
model satisfy covariance: it holds after scalar extension, and
extension is injective on the coordinates. The cyclic module
\[
 C=\mathcal T x_0\subset M
\]
is finite free over \(O\), and its
generic extension lies in the selected Hecke-stable span.
Write \(z=[p^a x_0]\) in \(M/M_{\rm W}\). If \(z\ne0\),
choose a maximal ideal \(\mathfrak m\) containing its annihilator.
Then \((\mathcal T z)_{\mathfrak m}\ne0\).
The kernel of \(\mathcal T\to\operatorname{End}_O C\)
kills \(x_0\) and hence \(z\). Thus the inclusions of ideals
\[
     \ker(\mathcal T\to\operatorname{End}_O C)
                   \subseteq\operatorname{Ann}_{\mathcal T}(z)
                   \subseteq\mathfrak m
\]
identify \(\mathfrak m\) with a maximal ideal of the
finite image algebra \(A\).

In a splitting field simultaneously triangularize the action on
\(C\). The diagonal characters \(\chi_j\) are integral,
and the matrices of \(J=\bigcap_j\ker\chi_j\) are strictly upper
triangular. In a triangular basis, with \(C_i\) the span of the first
\(i\) vectors, this gives
\[
 JC_i\subseteq C_{i-1}\quad(1\le i\le\dim C),
 \qquad J^{\dim C}C=0 .
\]
The action of the image algebra is faithful, so \(J\) is nilpotent. The product of integer rings
of the \(\chi_j\) is integral over \(A/J\). Lying over
gives a coordinate character with reduction \(\mathfrak m\).
Its generalized eigenspace is nonzero. On the generalized eigenspace \(V\), let \(a\) be the last exponent with \(\mathfrak n^aV\ne0\), for the maximal ideal \(\mathfrak n\) of the character. Then
\[
 0\ne\mathfrak n^aV\subseteq\ker\mathfrak n,\qquad \mathfrak n^{a+1}V=0.
\]
Any nonzero vector in this term is an eigenvector. It is an eigencharacter in the selected span:
extension of the interpolation frame preserves independence
because its nonzero full-rank minors remain nonzero.

\proofstep{Descent of the Weyl value test to the given row.}
Use the local models at all maximal ideals arising in this
way, equivalently those under integral eigencharacters occurring
in the selected span. If \(p^a\) gives the Weyl inclusion in
every such model, a maximal ideal detecting a nonzero class of
\(p^a x_0\) would contradict that inclusion. Thus \(p^a x_0\) has
Weyl values. Descent of the value tests gives the inclusion on
\(\mathcal C,\mathcal D\). In a localized coefficient basis the
condition is precisely nonnegative valuation on each coordinate.
It follows at the original level by the inverse-image description
above. This argument requires the same value and operator readings on both lattices.
\label{cor:lift-exponent-models-2}

\end{proof}

\paragraph{The two weights seen at a maximal ideal.}
At the chosen label put \(r=k_c+j+t+q\) and number the
coordinates from zero. Let \(A\) be the coefficient algebra
generated by the level matrices and the individual row-weight
projectors in the tensor. For a simple factor \(L\) of the
reduction, take decreasing real numbers \(\xi_{\tau,a}\) at
each label \(\tau\) and choose a joint weight \(\mu\) maximizing
\(\sum_{\tau,a}\xi_{\tau,a}\mu_{\tau,a}\) on \(L\).

If
\(\mu_{\tau,a}<\mu_{\tau,b}\) for \(a<b\), the transposition
of rows changes this sum by
\[
        (\xi_{\tau,a}-\xi_{\tau,b})
                     (\mu_{\tau,b}-\mu_{\tau,a})>0 .
\]
The row permutation preserves the weight set, so the
maximizer is dominant. The highest-weight theorem for this
coefficient algebra embeds \(L\) into the reduced dual Weyl
tensor of weight \(\mu\) and preserves the projectors. At one
label the theorem and the weight bounds give
\[
 \sum_a\mu_a=\sum_a\lambda_a,\qquad
 \lambda_{r-1}\le\mu_a\le\lambda_0,\qquad
 U_a=\lambda_a+r-1-a,\quad U_a'=\mu_a+r-1-a .
\]
In particular \(U,U'\) are strictly decreasing, of equal sum,
with entries in \([0,h_0]\), using \(\lambda_0+r-1\le h_0\).
\label{cor:lift-crystalline-depth-1}

Use the central Capelli generating polynomial
\(C_\nu(X)=\prod_{a=0}^{r-1}(X+\nu_a-a)\).
Its polarizations are coefficients of elementary row additions
followed by projection to the changed weights. The common degree
bound \(|\lambda|+|\mu|\) lets one use the same coefficient
operators on the reduced source and image. On a nonzero image
of \(L\) the two Capelli actions therefore have the same scalar
at every coefficient. The scalars are integral, so this gives
\[
       C_\lambda(X+r-1)\equiv C_\mu(X+r-1)\pmod p,
       \qquad \prod_a(X+U_a)=\prod_a(X+U_a')\ \text{in }k[X].
\]
Use the residue field \(k\) of the coefficient DVR here. Equality
of the two products gives the same multiset of residues, including
when several labels have the same residue.
\label{cor:lift-crystalline-depth-2}

Suppose the forms with values in \(L\) survive at the residual
ideal \(\mathfrak m\). At a chosen double coset \(c\), lift a
reduced value to the integral \(\mu\)-lattice. Its stabilizer
\(H_c\) has order invertible in \(O\). Averaging with the
level action \(R_c\) gives the lift
\[
        v_c^{\rm inv}=|H_c|^{-1}\sum_{h\in H_c} R_c(h)v_c .
\]
The reduction of \(v_c^{\rm inv}\) is the required value, since
the reduced value satisfies the stabilizer equations. Extend these
lifts by covariance. Let \(M_\mu\) be the integral weight-\(\mu\) forms and \(\bar v\) the surviving residual eigenform. Commutation of Hecke translations with reduction gives
\[
 \operatorname{Ann}(M_\mu)\subseteq\operatorname{Ann}(\bar v)=\mathfrak m.
\]
The finite torsion-free Hecke algebra argument
above now gives an integral characteristic-zero eigensystem
over \(\mathfrak m\), realized by a weight-\(\mu\) form since the
coordinate maps preserve independence. A residue-field embedding
aligning the two residual systems lifts to a local Galois
automorphism. This gives the congruence on every generator.
\label{cor:lift-crystalline-depth-3}

\paragraph{Automorphic and crystalline realizations at the chosen place.}
Use the coefficient field and the twisted tuple of this
lift at \(w\mid p\). Choose a \(p\)-adic embedding \(j_w\) inducing \(w\) and an isomorphism \(\iota:\overline{\Q}_p\to\C\). On the coefficient field their composite is prescribed:
\[
 \iota\circ j_w=j_{\mathbb C},
\]
where \(j_{\mathbb C}\) is the selected complex embedding. The latter comparison is algebraic.
The finite Hecke operators can be read through it because
their matrices on the finite class model are formed by
the same right translations and transported level
matrices.

At each member of the selected CM half, use
the finite adelic reading extending that \(p\)-adic
embedding. In the adapted basis the level matrices and
their inverses have integral entries. The twist is a
finite character, so adjoining the roots of unity of its
order puts these matrices and character values in the
same ring of integers \(O\). Evaluation at the finite
double classes gives the common integral form model.
The selected tame level makes the stabilizer orders units
in \(O\), including for the smaller weight \(\mu\).
\label{cor:lift-frame-galois-1}

\begin{proposition}[The oriented crystalline frame and the marked string]
\label{prop:lift-frame-crystalline-reading}
In the selected definite model, with the prescribed algebraic
type and spherical local factors over \(p\), the attached
crystalline representation has the rank-\(r\) normalized
Satake readings. For the \(\lambda\) system assume the specified
length-\(t\) divisor of the Hecke polynomial, \(r<2t\), and the
isolated length-\(t\) run in the weight at every label. Its
unique long Speh factor gives the marked consecutive crystalline
string. After the common character twist the \(\mu\)-frame has the following jumps at the two labels, in the given orientation:
\[
 \text{jump at }\tau=h_0-U'_{\tau,\bar a},\qquad\text{jump at }c\tau=U'_{\tau,a}.
\]
The marked string uses this orientation and the Frobenius transport of the proof.
\end{proposition}

\begin{proof}
\proofstep{The Hecke action after transport to the definite group.}
The characteristic-zero coefficient model first gives
a vector-valued form on the original definite space.
The isometry to the standard definite group transports
its rational and level actions. Use the algebraic
coefficient action at infinity to extend the finite
class values. Pair them using the positive invariant
Hermitian form on the coefficient type and take a
nonzero irreducible component of the resulting weight
space, using an invariant orthogonal complement.
The transported finite Hecke sums are the same on these
components.

In the split reading the prime is the
conjugate of the marked row prime. An isometry of the
two split groups identifies the right coset sums using
the dual double-coset label, which sends the list of elementary-divisor exponents to
\[
 (a_1,\ldots,a_r)\longmapsto(-a_r,\ldots,-a_1).
\]
This reads the given
Hecke eigenvalue, under \(\iota\), at the conjugate
prime. At primes of \(E\) above \(p\) the selected
level is spherical and the level character is trivial
on the corresponding hyperspecial group. The operators
indexed by good split primes away from \(p\) include
all coefficients of their Satake polynomials.
\label{cor:lift-frame-galois-2}

\proofstep{Attachment to the normalized Speh factors.}
Use unitary base change
with these local readings, the discrete spectrum
classification for \(\mathrm{GL}_r\), and Galois
attachment for the regular algebraic polarized cuspidal
constituents. Writing \(F=E^+\) for the totally real field of the definite group, the field hypotheses are
\[
 [F:\Q]\ge2,\qquad E/F\text{ unramified at all finite places}.
\] Base change has the
archimedean infinitesimal characters of the selected
algebraic type and the spherical components over \(p\).
Its discrete factors are Speh representations with
conjugate self-dual unitary cuspidal data. The spherical
component of the induced local product forces the
cuspidal local factors to be spherical.

For a cuspidal datum \(\pi\) of rank \(a\) in a Speh factor
of length \(d\), put
\[
 s_j=\frac{r-1}{2}+\frac{d-1}{2}-j,
 \qquad c_j=\frac{a-1}{2}-s_j \qquad (0\le j<d).
\]
In the unitary infinitesimal parameter let \(z\) range over
the exponents of \(\pi\) at a label. Regular algebraicity of
the full rank-\(r\) parameter gives
\[
 s_j-z\in\Z,
\]
because
\(z+(d-1)/2-j\) occurs there. The same regularity makes the
cuspidal exponents distinct at every label. The twist
\(\pi_j=\pi\otimes |\det|^{c_j}\) has exponents \(z+c_j\)
in \((a-1)/2+\Z\), and is polarized by the norm-power
character descending from \(|\cdot|^{2c_j}\). If \(a\) is odd,
multiply the polarizing character on the real subfield by
the quadratic character of the CM extension to give the
required sign at every real place. Its pullback is trivial.

Apply
Barnet-Lamb--Gee--Geraghty--Taylor, Theorem~2.1.1, to this
regular algebraic polarized cuspidal representation
in rank \(a\). With unitary Satake roots \(\alpha_{\pi,h}\) and Hodge labels
\(z_{\pi,\tau,h}\), the reciprocal normalization in rank \(a\)
reads
\[
 \operatorname{HT}_\tau(\rho_{\pi_j})
     =\{(a-1)/2-z_{\pi,\tau,h}-c_j\}_h
     =\{s_j-z_{\pi,\tau,h}\}_h,\qquad
 \operatorname{Sat}(\rho_{\pi_j})=\{q^{s_j}\alpha_{\pi,h}\}_h
\]
at the indicated good split prime for the Satake expression.
The first equality is the Hodge--Tate normalization and the
second multiset uses \((a-1)/2-c_j=s_j\).
The representation is crystalline above \(p\) because
\(\pi_j\) is spherical there. At a good split prime of
norm \(q\), the twist and the rank-\(a\) local Langlands
normalization multiply the unitary Satake roots by
\[
 q^{-c_j+(a-1)/2}=q^{s_j}.
\]

The direct sum over the cuspidal data and \(j\) therefore
has the full rank-\(r\) Satake roots scaled by
\(q^{(r-1)/2}\), with the corresponding full set of
crystalline jumps. Purity in Theorem~2.1.1 gives equal
absolute values for the unramified roots in each
cuspidal factor. If the cuspidal roots in unitary normalization have the common absolute value \(A>0\), then their product has modulus
\[
 A^a=1,
\]
where \(a\) is the rank of this cuspidal factor. This equality uses the unitary central character. Thus \(A=1\).
\cite{labesse-cm,blgg-weight}\label{cor:lift-frame-galois-3}

\proofstep{Identifying the unique long rank-one string.}
For the selected system of weight \(\lambda\) there is
additional information. Its occurrence in the chosen
block of the row supplies, at a good split prime, the
divisor of its Hecke polynomial coming from the
length-\(t\) string. Its center has modulus one in
unitary Satake normalization. To identify the global
factor, each cuspidal Satake multiset in the discrete
decomposition has unit modulus at this prime, by
temperedness of the cohomological conjugate self-dual
cuspidal factors in their unitary normalization. If a Speh constituent has length \(s\) and cuspidal rank \(a\),
its largest root modulus is \(q^{(s-1)/2}\), and its rank
contribution is \(as\). A root of modulus \(q^{(t-1)/2}\)
therefore requires
\[
       s\ge t,\qquad as\le r<2t .
\]
Hence \(a=1\), and two such constituents cannot fit in rank \(r\).

At each
archimedean label the specified length-\(t\) constant
run in the highest weight has strict drops at its two
ends when those ends are present. The shifted
infinitesimal exponents at this label are distinct.
The long Speh constituent gives a consecutive interval
of exponents. It cannot cross either drop, because each
drop leaves an absent intermediate exponent. The other intervals together have at most \(r-s<t\)
entries. They cannot cover the designated consecutive run of
\(t\) entries. The long interval must meet that run and cannot
cross a missing exponent at either boundary, so \(s\le t\).
Combining the size bound for the constituent and the two boundary gaps gives
\[
 t\le s\le t,\qquad s=t.
\]
Thus the long interval is exactly the specified run. Its Galois factors are the \(t\) successive
crystalline twists of a rank-one character. The
remaining Speh factors give \(\rho\).
\label{cor:lift-frame-galois-4}

\proofstep{Filtration jumps for the two CM orientations.}
The common normalization fixes their readings on both
orientations. Write \(\tau\) for a member of the
\(p\)-adic half chosen by the frame, \(c\tau\) for its
conjugate, \(s_\tau\) for the determinant weight shift
there, and
\[
 \bar a=r-1-a
\]
for a reversed coordinate.
Before the common twist, the type on \(\tau\) has highest
weight \(\mu_{\tau,a}+s_\tau\) and the type on \(c\tau\)
has its dual. The algebraic Hecke character with
crystalline jumps \(h_0-r+1+s_\tau\) on \(\tau\) and
\(-s_\tau\) on \(c\tau\) is chosen unramified over \(p\).
The sum of the two jumps is constant on all pairs, as
required for the algebraic character. Under the
reciprocal Frobenius convention the untwisted jump is
\(\bar a-\lambda_{{\rm aut},\bar a}\) in either local
highest weight \(\lambda_{\rm aut}\).

Adding the character
jump gives the two arrays
\[
 J_{\tau,a}=h_0-U'_{\tau,\bar a},\qquad
 J_{c\tau,a}=U'_{\tau,a},\qquad
 U'_{\tau,a}=\mu_{\tau,a}+r-1-a .
\]
Thus the marked orientation for \(U'\) is the conjugate
of the half used for the complex automorphic frame.
Apply the same twist to the selected \(\lambda\) system.
On that orientation the marked line jumps are
\(B_\tau+s\), \(0\le s<t\), with \(B_\tau\) the bottom
of the shifted run of \(\lambda\). On the other orientation
their bottom is \(h_0-B_\tau-(t-1)\), with the order of
the marked indices reversed. Both arrays and both
families of lines therefore lie in \([0,h_0]\).
The Frobenius permutation acts on all embeddings of the
unramified host and commutes with their CM involution.
This determines the transported label and its
orientation together in the crystalline comparison.
\label{cor:lift-frame-galois-5}
\end{proof}


\begin{lemma}[Identifying the marked factors in the companion reduction]
\label{lem:companion-residual-factors}
On the common coefficient extension of the two normalized
Hecke systems, the residual semisimplification is
\[
 \bar r_\mu^{\rm ss}\simeq \bar\rho^{\rm ss}\oplus\bigoplus_{s=0}^{t-1}\bar\chi_s.
\]
One may keep the marked occurrence of each scalar line in a
composition series of its actual reduced lattice.
\end{lemma}

\begin{proof}
At the good degree-one Frobenius classes, congruence of
the two Hecke systems gives
\[
 \det(X-r_\mu(g))\equiv
 \det(X-\rho(g))\prod_{s=0}^{t-1}(X-\chi_s(g))
                 \pmod {\mathfrak m_{\overline{\Q}_p}} .
\]
Indeed the reversed polynomial has Hecke coefficients times
signed integral norm powers and powers of the common twisting
character, which has unit values by compactness. For degree \(r\), reversing by
\[
 P(X)\mapsto X^rP(X^{-1})
\]
is a permutation of the coefficient list. It therefore
preserves the congruence after the common unit normalizations. The congruence
locus is closed because all coefficients are continuous and
the open unit ball in an ultrametric field is closed.
Degree-one Chebotarev density after finitely many omissions
extends it to every group element, also for the reciprocal
Frobenius convention. Choose a common finite coefficient field
and stable lattices.

For the two residual representations \(\bar R_1,\bar R_2\), Brauer--Nesbitt gives
\[
 \det(1-X\bar R_1(g))=\det(1-X\bar R_2(g))\ (\forall g)\ \Longrightarrow\ \bar R_1^{\rm ss}\simeq\bar R_2^{\rm ss}.
\]
A composition series therefore has the
prescribed line characters among its factors with their
multiplicities. Keep their positions in the filtration.
The block comparison takes these lattices over a common
finite coefficient extension with the marked characters
defined there. Compactness gives an invariant lattice
by taking the span of translates of one integral basis.
The coefficient field can be enlarged to contain the
labeled fields and the primitive \(p\)-th roots used in
the cyclotomic frame. This computation uses the original
residual filtration on the \(r_\mu\) side, with its
marked factors identified by the semisimplification.
\label{cor:lift-crystalline-depth-4}
\end{proof}

\paragraph{What the residual lines imply about the head.}
\hypertarget{source-lift-weight-exchange}{}
Let \(H=h_0\) and \(K=\lfloor H/(p-1)\rfloor\).
Here \(p>y>8\), \(p\) is unramified in the CM field, and
the large-prime inequalities give \(t+2K<p\).
Use all embeddings of the CM field at \(p\). At the conjugate embedding use the decreasing sequence obtained by reflection and reversal. In zero-based coordinates the operation is
\[
 (u^c)_a=H-u_{r-1-a}\quad(0\le a<r),\qquad (u^c)^c=u .
\] The
bottom of the new block is \(H-(B+t-1)\). Write
\(\sigma i\) for the source label of the Frobenius matrix
ending at \(i\).

In the local cyclic indexing it is \(i-1\).
This permutation commutes with conjugation. The Galois attachment can use either
orientation for the whole family. Pulling back the two
sequences and the block along that involution identifies its
\(J\) with \(U'\), and its line weights with
\(b_{i,s}=B_i+s\) at each embedding \(i\). A calculation at
the transported label then reads at \(i,\sigma i\) in the
original labels. Reflection preserves the characteristic-polynomial comparison. Indeed, for any list of residues of length \(r\),
\[
 \prod_{a=0}^{r-1}\bigl(X-(H-u_a)\bigr)=(-1)^r\prod_{a=0}^{r-1}\bigl((H-X)-u_a\bigr).
\]

\begin{proposition}[A depth bound from the residual lines]\label{prop:residual-head-depth}
For the shifted highest weights, conjugate labels and consecutive
line block just specified, the deficit at a head cut satisfies
\[
 \sum_{a<k_c+j}(\lambda_{i,a}-\mu_{i,a})
   \le qH+(p m_{\sigma i}+m_i)(t+2K).
\]
Here \(m_i\) counts the original shifted weights outside the consecutive
block at \(i\) whose residues occur in the block enlarged by \(K\)
at each end. The count \(m_{\sigma i}\) uses the Frobenius source label. At the indices satisfying the shape and
prime exclusions of Section~\ref{subsec:lift-exponent-primes},
both corrections vanish and the depth is at most \(\lfloor qh\rfloor\),
uniformly over the occurring Hecke ideal and the weight \(\mu\).
\end{proposition}

\begin{proof}
Use the norm coordinate of Section~\ref{subsec:cut-wach-models},
written \(X=Z\) in this calculation, with the same terminal
Hodge labels. Apply its construction to
the lattice of the congruent weight \(\mu\), to the inverse
images of the residual flag and to the specified scalar lines.
Their coefficient extension is over the selected place. Their
jumps are respectively \(U'_{i,a}\) and \(b_{i,s}\) and lie
in \([0,H]\). Write \(F\) for the residue field. In the resulting
\(F[[X]]\) lattice the reduced Frobenius has an integral right
inverse up to \(X^H\). For any \(d\)-rowed minor \(M\) of the reduced Frobenius matrix the elementary divisors give
\[
 \operatorname{ord}_X M\ \ge\ s_i(d),\qquad s_i(d)=\sum_{a=r-d}^{r-1}U'_{i,a}.
\]
The first cyclotomic matrix coefficient \(\Theta_i\) has
characteristic roots \(-\bar c U'_{i,a}\).
The residual filtration gives the same ranks and marked lines in
\(F((X))\) by exact 'etale realization.

Its intersection with
\(F[[X]]^r\) is a saturated flag, so extending bases at successive
steps gives one basis adapted to it. Each marked line occupies its own
diagonal block. Its two chosen generators in the Laurent
line differ by a unit times \(X^{\kappa_{i,s}}\), with the
sign fixed so that Frobenius in the adapted basis has order
\[
        \alpha_{i,s}=b_{i,s}+p\kappa_{\sigma i,s}-\kappa_{i,s}.
\]
In fact Frobenius sends the parameter to \(X^p\), and the
reduced rank-one Frobenius is a unit times \(X^{b_{i,s}}\).
The latter follows from the rank-one determinant, whose
order before reduction is \(b_{i,s}\) at \(X-p\) by height
and the local elementary divisor. Integrality gives
\(\alpha_{i,s}\ge0\) and the height inverse gives
\(\alpha_{i,s}\le H\) on the singleton block. The inverse
is block triangular after inverting \(X\), so the two
diagonal entries multiply to \(X^H\).

\proofstep{From the determinant to the minor estimate.}
To constrain individual head weights we pass from the determinant to a minor.

For the cyclotomic generator, \(\gamma(X)/X\) has first
terms \(1+\bar c X\) modulo \(X^2\) in this reduced
coordinate. For a unit \(v=v_0+v_1X+O(X^2)\), the constant and linear terms cancel in the quotient:
\[
 \gamma(v)\equiv v_0+v_1X\equiv v\pmod{X^2},\qquad \gamma(v)/v\equiv1\pmod{X^2}.
\] The coefficient of \(X\) on the same marked line
is therefore \(-\bar c(b_{i,s}-\kappa_{i,s})\).
An integral basis change conjugates this first coefficient by
its constant matrix, since the cyclotomic action fixes the
change of basis modulo \(X^2\).

Taking characteristic
polynomials of the block triangular matrix and cancelling
the nonzero common scalar shows that the residues
\(b_{i,s}-\kappa_{i,s}\), counted with multiplicity, occur
among those of \(U'_i\). The determinant of the Frobenius
minor on the marked positions is the product of their
diagonal entries. The minor divisibility then gives
\begin{equation}\label{eq:lift-line-minor}
       s_i(t)\le\sum_s b_{i,s}
              +p\sum_s\kappa_{\sigma i,s}-\sum_s\kappa_{i,s}.
\end{equation}
The flag and the rank-one comparison here come from the
same residual representation at the same maximal ideal.
\label{cor:lift-weight-exchange-1}

Since \(b_{i,s}\) and \(\alpha_{i,s}\) belong to \([0,H]\),
we have \(|p\kappa_{\sigma i,s}-\kappa_{i,s}|\le H\).
For fixed \(s\), let \(j\) maximize \(\lvert\kappa_{j,s}\rvert\).
Apply the relation at \(i=\sigma^{-1}j\), where the term
with coefficient \(p\) is \(\kappa_{j,s}\). The triangle inequality at this predecessor gives
\[
 p|\kappa_{j,s}|\le H+|\kappa_{\sigma^{-1}j,s}|
            \le H+|\kappa_{j,s}| .
\]
Hence every \(|\kappa_{i,s}|\le H/(p-1)\).
To bound the sum of the shifts, enlarge the block
\(I=[B_i,B_i+t-1]\) by \(K\) at both ends to \(I^+\).

Let \(m_i\) count original weights outside \(I\) whose
residue occurs in \(I^+\). The strict length bound makes
reduction modulo \(p\) injective on the integers in \(I^+\).
The \(t\) shifted entries \(b_{i,s}-\kappa_{i,s}\) lie
there. Match their residues to distinct occurrences among
the \(U_i\) using the two multiset comparisons. At most
\(m_i\) are matched outside \(I\). All other matches are
equalities in \(I^+\). Pair the remaining shifted entries
arbitrarily with the unused block entries. Each difference
is bounded by \(t+2K\), so
\[
                \left|\sum_s\kappa_{i,s}\right|
                            \le m_i(t+2K).
\]
This also holds at \(\sigma i\), where the exceptional count
uses its own block. Reflection preserves the count because it
carries the block, its enlargement and residue congruence
bijectively to their reflected versions.

\proofstep{Returning to the head of the weight.}
The deficit of the first \(k_c+j\) weights equals the excess
of the rest by equality of total sums. Of the last \(t+q\)
entries of \(U'_i\), the smallest \(t\) have sum \(s_i(t)\)
and the other \(q\) are at most \(H\) apiece. In \(U_i\)
the corresponding total is at least \(\sum_s b_{i,s}\)
because the bottom \(q\) entries are nonnegative. Equation~\eqref{eq:lift-line-minor}
thus proves
\[
 \sum_{a<k_c+j}(\lambda_{i,a}-\mu_{i,a})
   \le qH+(p m_{\sigma i}+m_i)(t+2K).
\]
The staircase cancels term by term on the left. The shape
and prime exclusions of Section~\ref{subsec:lift-exponent-primes}
give \(m_i=m_{\sigma i}=0\). Every surviving highest
weight at this index with a larger weight therefore has depth at most
\(i_i=\lfloor qh\rfloor\), with the choice of
\(\mu\) and of the occurring Hecke ideal unrestricted.
\label{cor:lift-weight-exchange-2}
\end{proof}


\paragraph{The localized tensor inclusion.}
For each embedding \(\tau\), let \(V_\tau=\nabla(\lambda_\tau)\)
be the dual Weyl module of highest weight \(\lambda_\tau\) over \(O\).
Write \(B_\tau=\Delta(\lambda_\tau)\subset V_\tau\) for its Weyl lattice.
Write \(V=\bigotimes_\tau V_\tau\).

\paragraph{Divisibility on one row-weight space.}
\hypertarget{source-lift-pieri-cost}{}
We give the bounded-depth computation over \(\Z_{(p)}\).
For a dominant weight \(\nu\), a row weight \(v\) and a vector
of coefficients \(\xi=(\xi_a)_a\), put
\(d_\xi(v)=\sum_a\xi_a(\nu_a-v_a)\).
In the application, \(\xi_a=1\) on the first \(k_c+j\) coordinates
and \(\xi_a=0\) on the others.

\begin{lemma}[Rectangular Pieri denominator]
\label{lem:rectangular-pieri-denominator}
Consider first a rectangular
initial block \(\nu=(w^b,\beta)\), where \(\beta\) is a
partition of length at most \(d\) and \(\beta_1\le w\).
An upper bound for the extra exponent in passing from the
last \(d\) coordinates to the full coefficient module at
\(n=bw-\sum_{a<b}v_a\) is
\begin{equation}\label{eq:lift-pieri-polynomial}
 v_p(D_n(w)),\qquad
 D_n(Z)=\prod_{c}(Z-c)^{m(c)} .
\end{equation}
Here one takes partitions \(\tau\supset\beta\) with
\(|\tau/\beta|=n\), column lengths of \(\tau/\beta\) at
most \(b\), and no added box of content \(w\). Content is
column minus row, both numbered from zero. The integer
\(m(c)\) is the maximum number of boxes of content \(c\)
in such a skew diagram, with empty maximum zero.

In
particular \(D_n(w)\ne0\). If \(p^e\) suffices on the
Levi weight spaces with \(d_\xi\le I\) (extend their weights
by the initial rectangle), then \(p^{e+v_p(D_n(w))}\)
suffices on the indicated full weight space whenever
\(d_\xi(v)\le I\) and \(\xi\) is decreasing in the weak sense.
\label{cor:lift-pieri-cost-1}

\end{lemma}

\begin{proof}
\proofstep{Integral raising and lowering tests.}
Use the polynomial representation over \(\Q\) with the
contravariant Fock form normalized to give the highest
weight vector norm one. The Levi highest subspace \(W\)
is generated from this vector by the last \(d\) rows. The
first block acts there by the determinant character and the
upper cross-block operators kill it. Let \(F_{u\alpha}\)
be the cross-block lowering operator from first row \(\alpha\)
to last row \(u\), and \(E_{\alpha u}\) its adjoint.
The lowering operators commute with each other. Their products
of order \(n\) applied to \(W\) span the depth \(n\) subspace.
Indeed reordering words in row operators puts the lowering
operators on the left and the raising operators on the right.
Let \(T\) be the sum of the first diagonal operators. On the span of depth \(n\) it acts by \(bw-n\). Since it is self-adjoint in the contravariant form, vectors at different depths satisfy
\[
 (n'-n)\langle v_n,v_{n'}\rangle=\langle Tv_n,v_{n'}\rangle-\langle v_n,Tv_{n'}\rangle=0 .
\]

For a multiset \(m\) of \(n\) pairs let
\(m!=\prod_{u,\alpha}m_{u\alpha}!\).
The lowering generators of the integral lattice are
\(F^m/m!\) on the Levi Weyl lattice. Each is a product of
integral divided row powers. By adjunction, the test against a divided lowering word is a test in the dual Levi lattice. For a Levi vector \(z\),
\[
 \left\langle v_0,\frac{F^m}{m!}z\right\rangle =\langle y_m,z\rangle,\qquad y_m=\frac{E^m}{m!}v_0.
\]
These entries determine \(v_0\), by the spanning assertion
and nondegeneracy. In the polynomial model, restriction to
the Levi block identifies its highest weight generator and
its divided row words with their smaller-rank counterparts.
The restricted dual tests are integral there because their pairings
with the transposed words read the coefficients on the unit
lower triangular chart. The ordered row additions extract
these coefficients and restriction is injective on \(W\).
\label{cor:lift-pieri-cost-2}

\proofstep{The commuting polynomial on multisets.}
To compute with the entries, represent a vector by an array
of ordered lowering words of length \(n\) with coefficients
in \(W\). Let \(L\) denote evaluation of the words and \(S\)
the map which sums the array over words of each multiset.
Thus \(L\) factors through \(S\). In tensor position \(j\),
numbered from one, set
\[
 J_j=\sum_{u,v}e^{(j)}_{uv}E^W_{vu}
                       +\sum_{a<j}(a\ j)_Y ,
 \qquad A(Z)=\prod_{j=1}^n(Z-J_j).
\]
Here \(e^{(j)}_{uv}\) acts on the last-block row index,
\(E^W_{vu}\) on the Levi coefficient and the transpositions
exchange just the last-block indices. The raising tests of
\(Lz\) satisfy
\[
                    \iota Lz=S A(w) z,\qquad
                    (\iota v_0)_m=E^m v_0/m! .
\]
For this identity commute the first raising operator past
each lowering in a word, using
\([E_{\alpha u},F_{v\beta}]
  =\delta_{uv}E_{\alpha\beta}-\delta_{\alpha\beta}E_{vu}\).

On the highest subspace the first term gives the scalar
\(w\), the second gives the Levi action. Commuting the
resulting row generators through the other letters gives
the transpositions in \(J_j\). Iteration over the letters
gives the ordered product on the sum over permutations of
the word. In the sum over all orders of a fixed multiset, the stabilizer permutes the repeated copies of each pair. Its cardinality is
\[
 |\operatorname{Stab}(m)|=\prod_{u,\alpha}m_{u\alpha}!=m!,
\]
giving the divided raising test above on collecting equal words.
These are identities over characteristic zero. The
coefficients of \(A(Z)\) descend to commuting operators on
multiset arrays. Indeed the successive partial quadratic
Casimirs commute, their differences give the \(J_j\), and
adjacent transpositions together generate the permutations.
Conjugating the defining sums by the adjacent transposition on the last-block indices gives the two relations
\[
 sJ_a=J_{a+1}s-1,\qquad sJ_{a+1}=J_as+1,\qquad s=(a\ a+1).
\]

The transposition \(s\) commutes with the other \(J_i\).
Using the two swap relations and commutativity of \(J_a,J_{a+1}\)
gives
\[
        s(J_a+J_{a+1})=(J_a+J_{a+1})s,\qquad
        sJ_aJ_{a+1}=J_aJ_{a+1}s .
\]
These are the two coefficients that show commutation with
\((Z-J_a)(Z-J_{a+1})\), hence with \(A(Z)\). Swaps on the
first-block indices commute with all the \(J_i\). Its coefficients preserve integral
arrays since the operator upstairs uses integer permutations
and Levi row operators and \(S\) is onto over the integers.

\proofstep{Contents on the surviving Pieri paths.}
Tensoring the irreducible Levi module successively with
the defining representation gives the rational Pieri paths
which add one box per step. On a path the eigenvalue of
\(J_j\) is the content of the box just added. In fact the
quadratic Casimir on shape \(\kappa\) has scalar
\(\sum_a \kappa_a(\kappa_a+d-1-2a)\). For the box in row \(a\) the new column has zero-based index \(\kappa_a\).
The change of the form, minus \(d\), is
\[
 (\kappa_a+1)(\kappa_a+d-2a)
       -\kappa_a(\kappa_a+d-1-2a)-d=2(\kappa_a-a).
\]
Here the content is \(\kappa_a-a\), equivalently the
column number \(\kappa_a+1\) in one-based column notation
minus one, then minus \(a\). The tensor Casimir identity gives twice \(J_j\)
for this operator difference.

To see which paths can survive
\(S\), fix the counts \(n_\alpha\) of the first-block letters.
The multiset arrays with these counts, as a module of the
last block in characteristic zero, are
\(W\otimes\bigotimes_{\alpha=1}^b
\operatorname{Sym}^{n_\alpha}(\Q^d)\).
Indeed their multiset basis is indexed by one monomial per
first-block letter. By the horizontal-strip Pieri rule, let \(\lambda^{(j)}/\lambda^{(j-1)}\) be the strip contributed by the \(j\)-th symmetric power.

In each column \(c\),
\[
 (\tau/\beta)'_c=\sum_{j=1}^b\bigl((\lambda^{(j)})'_c-(\lambda^{(j-1)})'_c\bigr)\le b,
\]
since each summand is zero or one. Here the left side counts the boxes of the skew diagram in that column. The column bound kills the image under \(S\) of every
\(\tau\) with a column of length \(>b\) in \(\tau/\beta\).
For a surviving path, each successive \(J_i\) reads the
content of its added box, so the scalar on its \(\tau\)-space is
\[
       P_\tau(Z)=\prod_{\square\in\tau/\beta}
                   (Z-\operatorname{cont}(\square)).
\]
The product depends only on the skew shape. Since \(S\)
is equivariant, \(S A(w)\) has image in the sum of the
surviving spaces with \(P_\tau(w)\ne0\).
\label{cor:lift-pieri-cost-3}

\proofstep{The common content polynomial clears the inverse.}
There is an integral way to invert this last operation up
to the scalar \(D_n(w)\). Divide \(D_n(Z)\) by the universal
monic polynomial \(Z^n+\sum_{a<n}c_a Z^a\) over
\(\Z[c_0,\ldots,c_{n-1}]\), and evaluate the quotient
at \(w\) with the commuting coefficient operators of \(A\)
in place of \(c_a\). Call the resulting integral operator
\(Q\).

On every path occurring in the image,
\(P_\tau\mid D_n\) by the content maxima. The specialized
remainder is zero. The polynomial-division identity on a surviving path thus specializes to
\[
 A(w)Q=D_n(w)\,\mathrm{id}
\]
on the image under consideration.
This division incurs no denominators for spectral
projectors. The coordinate of a raising test has weight
\(v+\sum m_{u\alpha}(e_\alpha-e_u)\).
Each \(J_j\) respects the matching between coordinate
indices and their required weights, since its matrix unit
and Levi action have opposite weight changes. The depth condition is preserved since \(\xi_\alpha\ge\xi_u\) at each contributing pair. The weight shift in a raising coordinate has pairing
\[
 \left\langle\xi,\sum_{u,\alpha}m_{u\alpha}(e_\alpha-e_u)\right\rangle=\sum_{u,\alpha}m_{u\alpha}(\xi_\alpha-\xi_u)\ge0 .
\]
Thus under \(Q\) each coordinate remains in the dual Levi lattice at the allowed depth.
Multiplication by \(p^e\) puts these coordinates in the
Levi Weyl lattice.

Apply the lowering-word evaluation
to the multiset array \(p^e Q\iota v_0\). It belongs to
the full Weyl lattice because each undivided word is an
integer multiple of its divided word. Its raising tests
are \(p^e D_n(w)\iota v_0\), so injectivity of the tests
proves \eqref{eq:lift-pieri-polynomial}. The scalar prime
to \(p\) is a unit. The descent to the smaller rank here
uses the indicated polynomial Weyl lattices with their
normalized pairing.\label{cor:lift-pieri-cost-4}

\end{proof}

\paragraph{Counting the factors at the selected weight.}
\begin{lemma}[Depth bounds at the selected rectangular block]
\label{lem:selected-pieri-depth}
On the selected row-depth cutoff \(I=i_i\), write \(c_{\rm out}\) for the cost of the two parts outside the rectangle. Then
\[
 c_{\rm out}\le\frac{(j+q)(h_0-1)}{p-1}.
\] The rectangular factors are counted by
the shifted intervals below, for each integer \(0\le n\le I\).
In the nonexceptional range \(h_0<p^2\) this gives the stated
coarse exponent for the projected inclusion, using the original
Frobenius label and the same grading.
\end{lemma}

\begin{proof}
\proofstep{Complementary coordinates and the one-row costs.}
For a rectangular block at the end reverse the coordinates
and take complementary weights about a common bound \(M\).
This is contragredience with a determinant twist. It
interchanges the two block positions, replacing \(\xi\) by
its negative reverse and preserving the graded deficit.
In a one-row step the added skew diagram is a horizontal
strip. Two boxes of the same content, farther right in a
lower row, would force a common added column. Thus all
contents in this strip are distinct. The surviving boxes cannot reach content \(w\)
since each row starts at content at most \(w\) and must
avoid \(w\).

Taking the maximum of the shifted widths, the surviving factors therefore satisfy
\[
 1\le w-c\le h_0-1.
\] More precisely
for an appended row \(a\), numbered from zero, the
complementary rectangle is \(M-\lambda_a\), and the
smallest possible content is \(1-a\) for \(a\ge1\).
With \(M=\lambda_0\), all factors are at most
\(\lambda_0+a-\lambda_a-1\le h_0-1\).
In each such step the distinct positive integers \(w-c\)
lie in \([1,h_0-1]\), so their product divides \((h_0-1)!\).
The empty Levi contributes the product one. It follows that
\[
  v_p\!\left(\prod_c(w-c)\right)
  \le v_p((h_0-1)!)
  \le\left\lfloor\frac{h_0-1}{p-1}\right\rfloor .
\]
\label{cor:lift-pieri-cost-5}

Start with the full exponent of the first \(k_c\) rows
(zero for the selected nonexceptional prime). Append the next
\(j\) rows one at a time, using the zero grading. Append
the rectangle of length \(t\) at once, with
\(w=\lambda_{k_c+j}\) before reversal. For this step the
block depth equals the deficit above the rectangle by
equality of the two total weights. Finally append the \(q\)
bottom rows individually, retaining the grading by the
first \(k_c+j\) entries. The preceding argument with
decreasing \(\xi\) keeps the same allowance \(I=i_i\)
through these last steps. The middle step must therefore be bounded at precisely the integer depths
\[
 n\in\{0,\ldots,I\},\qquad I=i_i .
\] The two outside
steps together cost at most \((j+q)(h_0-1)/(p-1)\).

\proofstep{Which labels can occur at bounded depth.}
On the reversed Levi let \(M=\lambda_0\) and
\(\beta_b=M-\lambda_{k_c+j-1-b}\), with rectangle height
\(M-w\). For a box of content \(c\), put
\(x=M+t-1-c+q\). Subtract the top weight \(B_\tau=w+t-1+q\) of the original shifted block. The reversed height then contributes the same factor:
\[
 x-B_\tau=(M+t-1-c+q)-(w+t-1+q)=(M-w)-c .
\] The labels in a row start
one below its shifted weight and decrease by one at each
box. They are strictly above \(B_\tau\) by the omitted
content condition.

The first \(k_c\) original rows have
shifted weights in \([x_-,x_+]\) and the next \(j\)
at most \(x_-\). The latter rows supply at most \(j\)
boxes of each label below \(x_-\), none above or at that
endpoint. The first \(k_c\) rows supply at most \(k_c\) boxes of a
label, and at most \(t\) since boxes on one content
diagonal force a common added column. For \(x<x_-\) each
of these rows reaching \(x\) has at least \(x_--x\) added
boxes. Their total is at most \(n\). If \(c_x\) counts these rows reaching \(x<x_-\), then
\[
       c_x(x_--x)\le n\le I,\qquad
       c_x\le \min\{k_c,t,I/(x_--x)\}.
\]
This multiplicity estimate applies at each such label.

\proofstep{Counting multiples of the residue prime.}
The \(j\)-row contribution is bounded by
\(jv_p((x_--B_\tau-1)!)\le j(h_0-1)/(p-1)\).
All relevant positive \(x-B_\tau\le h_0<p^2\) have
valuation at most one. In \([x_-,x_+]\) the number
of multiples of \(p\) is at most \((x_+-x_-)/p+I_p\),
with the indicator from Section~\ref{subsec:lift-exponent-primes}.
For \(x<x_-\), use the representative of the residue of \(x_--B_\tau\) in \(\{1,\ldots,p\}\).

The positive residue representative is
\[
 r_p=1+\bigl((x_--B_\tau-1)\bmod p\bigr).
\] The
possible distances \(x_--x\) are \(r_p+ap\), strictly
less than \(x_--B_\tau\). The term \(a=0\) gives
\(\min(k_c,t,I/r_p)\) if this strict inequality holds.
For \(a\ge1\) sum \(I/(ap)\) to get at most
\((I/p)(1+\log(h/p))\) when \(p\le h\).
Enlarging the quantities in this bound gives the claimed coarse estimate:
\[
 k_c\le k,\qquad I\le qh,\qquad \frac{h_0-1}{p-1}\le\frac{h+r}{p-1}+6 .
\] The cost is an
integer, so one may take the floor. Finally the projected
inclusion over \(O\) follows from the integral standard
bideterminant spanning theorem. Each generator is homogeneous
in the row degrees, and its projection is either itself or
zero. Its prime-to-\(p\) clearing integer from the integral
computation becomes a unit in \(O\).
\label{cor:lift-pieri-cost-6}
\end{proof}


\begin{lemma}[Descent from projected Weyl values]
\label{lem:localized-projected-weyl-values}
With the integral weights of the tensor \(\bigotimes_\tau
\nabla(\lambda_\tau)\), use the projected inclusions of
Lemmas~\ref{lem:rectangular-pieri-denominator}
and~\ref{lem:selected-pieri-depth}. The Weyl inclusion on the actual forms at
a selected residual ideal is obtained by clearing each factor
by its full exponent, or by the smaller allowable depth cost
at a larger-weight factor. Write \(d_\tau\) for the charge so chosen at each label, \(V\) for the unrestricted tensor lattice and \(D_\varnothing\) for the lattice using the Weyl factor at every label. After passing to the forms at the selected residual ideal \(\mathfrak m\), the inclusion is
\[
 p^{\sum_\tau d_\tau}\,\mathcal F(V)_{\mathfrak m}
 \subseteq \mathcal F(D_\varnothing)_{\mathfrak m},
\]
where \(\mathcal F(L)\) denotes the module of forms with values in \(L\).
\end{lemma}

\begin{proof}
For a set \(S\) of
unrestricted labels, let \(D_S\) be the image of the tensor
using \(V_\tau\) on \(S\) and \(B_\tau\) elsewhere. Each
\(D_S\) is stable under the coefficient algebra and its forms
are stable under the Hecke translations. The exact exponents
\(e_\tau\) give \(p^{\sum e_\tau}V\subset D_\varnothing\)
by multiplication on pure tensors. Thus \(V/D_\varnothing\)
has finite length. Every simple factor of a subquotient of
this quotient occurs in the reduction of \(V\) modulo the
uniformizer \(\varpi\). Indeed, multiplication by a power of \(\varpi\) gives the surjection onto each successive quotient:
\[
 V/\varpi V\ \twoheadrightarrow\ \varpi^a V/\varpi^{a+1}V,\qquad [v]\longmapsto[\varpi^a v]\quad(a\ge0).
\]

Write \(P_{\tau,w}\) for the integral weight projector,
\(i_\tau\) for the depth allowance and \(c_\tau\) for the
integral bounded-depth cost just obtained. For
\(a\notin S\), set \(D=D_S\) and \(D'=D_{S\cup\{a\}}\).
If the full exponent is used at \(a\), then
\(p^{e_a}D'\subset D\) on pure tensors. At an index with a larger weight
where \(c_a<e_a\), form the actual intermediate lattice
\[
                          E=D+p^{c_a}D'.
\]
For any weight \(w\) at depth at most \(i_a\) we have
\(p^{c_a}P_{a,w}V_a\subset B_a\). Applying that projector
in just factor \(a\) shows \(P_{a,w} E\subset D\).

In a simple factor \(L\) of \(E/D\), the projector at its
highest weight \(\mu_a\) has nonzero image. This follows
from the highest-weight line in the coefficient classification.
If the highest weight had depth at most \(i_a\), the integral inclusion would give
\[
 P_{a,\mu_a}E\subseteq D\quad\Longrightarrow\quad P_{a,\mu_a}(E/D)=0,
\]
contradicting the nonzero image on \(L\). Thus all simple factors here have
excess depth. By the preceding crystalline estimate their
localized forms vanish. This argument requires stability of
the intermediate lattices and the integral action of the
projector on the full tensor.

Take a composition series for \(E/D\). At each successive
quotient map a form to the form obtained by projecting its
values to the simple factor. This map is Hecke-equivariant.
Filter \(E\) by its series \(D=D_0\subset\cdots\subset D_b=E\), and denote the projection at the last step by \(\pi_b\). The vanishing for the simple factor reads
\[
 (\pi_b\phi_b)_{\mathfrak m}=0,\qquad \pi_b:D_b\to D_b/D_{b-1}.
\]
There is \(s_b\notin\mathfrak m\) killing that image before
localization, so \(\phi_{b-1}=s_b\phi_b\) takes values in
\(D_{b-1}\). Composing the successive denominator clearings gives
\[
 (s_1\cdots s_b)\phi_b\ \text{with values in }D_0=D .
\]
Multiplying these denominators down the series proves the
localized inclusion of the forms with values in \(E\) into
those with values in \(D\).

In particular \(p^{c_a}\) moves
the forms on \(D'\) into the forms on \(D\) after localization.
Inducting over the unrestricted labels gives
\[
 p^{\sum_{a\ {\rm thick}}\min(e_a,c_a)+
       \sum_{a\ {\rm other}}e_a}
            M_{\mathfrak m}\ \subset\ (M_{\rm W})_{\mathfrak m}.
\]
All denominators used are units in this localization.
The simple-factor argument works separately at each integral
joint character over the ideal. The surviving bound is the
same for every such character in the prime range under
consideration. This proves the passage from the projected
coefficient inequalities to the lattice of actual forms.
\label{cor:lift-graded-inclusion-1}
\end{proof}



\begin{lemma}[Least exponent on the original Weyl value pair]
\label{lem:weyl-least-exponent-local}
For the coefficient models of Proposition~\ref{prop:local-weyl-value-inclusion},
let \(B_j(p)\) be the bound from Section~\ref{subsec:lift-exponent-primes}
for the weight record of lift \(j\). For the ramification index \(e_w\) and the least local exponent,
\[
 J_{j,w}/e_w\le B_j(p)\qquad(w\mid p).
\]
This normalized upper bound descends through the coefficient extension.
\end{lemma}

\begin{proof}
The exponent in these models is computed with the actual
Frobenius and conjugation labels of the frame. In the indicated
range the two window correction counts vanish by the shape
calculation of Section~\ref{subsec:lift-exponent-primes}.
The depth and all other terms are therefore independent of the Frobenius and conjugation labels. One power applies at every maximal ideal in the cover. Outside the weight range use the sum of the full Weyl exponents, with exponent zero above height. The small and excluded primes likewise use the full Weyl sum. At a remaining prime the contribution of each larger-weight index is the minimum of its full exponent and its integral coarse correction. These comparisons fix the
tuple at \(w\) before taking the least exponent and keep the
original weight record for the prime sum. In compatible Smith bases the two lattices to be compared are
\[
 \mathcal C_w=\bigoplus_i O_we_i,\qquad \mathcal D_w=\bigoplus_i O_w\varpi_w^{d_i}e_i\subseteq\mathcal C_w.
\]

Then
\[
     J_{j,w}=\max_i d_i,\qquad
     p^{B_j(p)}\mathcal C_w\subseteq\mathcal D_w
             \ \Longrightarrow\ d_i\le e_wB_j(p)\quad\text{for all }i.
\]
The same comparison applies if the surviving eigensystem
requires an extension of the coefficient field. Over the extension the exponents and
the absolute ramification index both multiply by its relative
ramification index.
\label{cor:lift-exponent-models-3}
\end{proof}

